\documentclass[preprint,12pt,authoryear]{elsarticle}
\usepackage[letterpaper,hmargin=0.9in,top=0.76in,bottom=0.78in]{geometry}
\usepackage[T1]{fontenc}
\usepackage{lmodern,textcomp}
\usepackage{booktabs,tabularx,array}
\usepackage{listings,needspace}
\usepackage[hidelinks]{hyperref}
\hypersetup{pdftitle={Supplementary Material: A Neuro-Symbolic Architecture for Safe and Feasible Self-Recovery in Multi-Agent Manufacturing Systems}}

\setlength{\parindent}{0pt}
\setlength{\parskip}{3pt}
\setlength{\intextsep}{6pt plus 1pt minus 1pt}
\raggedbottom
\urlstyle{tt}
\let\code\nolinkurl
\newcolumntype{Y}{>{\raggedright\arraybackslash}X}
\newcommand{\excerptlabel}[2]{\par\Needspace{4\baselineskip}\smallskip{\small\bfseries #1 excerpt\enspace{\normalfont\itshape #2}}\par\nopagebreak}
\lstset{
  basicstyle=\small\ttfamily,
  breaklines=true, breakatwhitespace=false, breakautoindent=false,
  columns=fullflexible, keepspaces=true, showstringspaces=false,
  upquote=true, xleftmargin=1em, aboveskip=2pt, belowskip=4pt
}
\lstdefinestyle{recorded}{
  basicstyle=\footnotesize\ttfamily, frame=leftline,
  framesep=5pt, framerule=0.35pt, breakindent=0pt,
  escapeinside={(*@}{@*)}
}
\newcommand{\recordheading}[1]{\par\Needspace{5\baselineskip}\medskip{\small\bfseries #1}\par\nopagebreak}
\newcommand{\papertitle}{A Neuro-Symbolic Architecture for Safe and Feasible Self-Recovery in Multi-Agent Manufacturing Systems}

\begin{document}
\pagestyle{plain}
\begin{center}
  {\Large\bfseries \papertitle\par}
  \medskip
  {\large Supplementary Material\par}
\end{center}

This supplement documents how the framework generates a recovery outline
after a manufacturing failure, composes the accepted events from resource
primitives, and executes the resulting programs in Gazebo. Appendices~A and~B
use Part slippage as a worked example, showing the prompts, selected model
replies, and the check findings that led to revisions. Appendix~C presents
the recovery events executed in all four failure scenarios.

The complete \hyperref[sec:outline-prompt]{outline prompt} starts on
page~\pageref{sec:outline-prompt}. The complete
\hyperref[sec:composition-prompt]{primitive composition prompt} for the gear
pick starts on page~\pageref{sec:composition-prompt}. Each includes the user
message and the system message that specifies the reply format.

\appendix
\section{Recovery Outline Synthesis}
\subsection{What the LLM receives}

The prompt gives the large language model (LLM) the observed failure, the recovery goal, and the recorded
states and capabilities of the resources and parts. The LLM can request further
facts before proposing events. Each method in Table~\ref{tab:methods} was tested
at both generation levels. All finished outlines were evaluated for feasibility,
safety and task continuation.

\begin{table}[h!]
\centering\small
\caption{Information available during outline synthesis.}\label{tab:methods}
\begin{tabularx}{\linewidth}{@{}l l Y@{}}
\toprule
Method & Safety rules in prompt & Checks and feedback during synthesis \\
\midrule
Baseline & No & None. \\
Informed baseline & Yes & None. \\
Framework & No & The Resource Agent (RA) checks feasibility. The Central Controller Agent (CCA) checks safety and nonblocking. The next prompt carries refusal findings. \\
\bottomrule
\end{tabularx}
\end{table}

\excerptlabel{Prompt}{Framework, event-level}
{\small\leftskip=1em\ttfamily Propose 2{\rmfamily\textendash}5 different ways to do the next step of that plan, most preferred first.\par}
The framework accepts the first candidate that passes the checks. The worked
example uses this setting with Claude Opus 5.5 at medium reasoning effort.

\excerptlabel{Prompt}{Framework, outline-level}
\begin{lstlisting}
Give the whole plan in `events`: every recovery event, in the order they happen.
\end{lstlisting}
Here the LLM proposes the entire sequence in one reply. If the framework
refuses it, the LLM receives the findings before revising the outline.

\clearpage
\subsection{Part slippage and the accepted outline}\label{sec:outline-case}

\code{ur5e-4} dropped \code{gear_small} while approaching the assembly board.
The turn~1 prompt located the loose gear and identified its intended slot.

\excerptlabel{Prompt}{Turn 1, failure observation and recovery goal}
\begin{lstlisting}
- gear_small: not held, not at any station, measured at (0.250, 0.450, 1.025)
\end{lstlisting}
\begin{lstlisting}
- gear_small: part_location assembly_board-v1 (target Gear_Plate/Gear_Shaft_1) (normally ur5e-4's place_insert; ur5e-4 failed)
\end{lstlisting}

The goal identified the slot, but the failure observation alone did not establish
which robot could reach the loose gear. On turn~1 the LLM requested eight facts,
including the gear's position relative to the robots.

\excerptlabel{Reply}{Turn 1, first two references in the context request}
\begin{lstlisting}
"decision": "need_context",
"context_requests": [
  "/parts/gear_small/pose_relations",
  "/parts/gear_small/state_variables",
\end{lstlisting}

The turn~2 prompt reported a horizontal distance of 0.255\,m from
\code{ur5e-3} and 0.982\,m from \code{ur5e-4}. Each listed a radius of
0.886\,m and a tolerance of 0.01\,m. The gear lay inside the first robot's
work envelope. That robot was already holding \code{KET4_Square_4mm}.

The LLM's first candidate was to place the square peg on the board. The RA
found the placement feasible. The CCA refused it under
\code{SAFE_shared_area_mutex} because \code{ur5e-4} still occupied the board
area. The second candidate staged the peg in the empty buffer zone and passed
both checks. It became event~1 in Table~\ref{tab:events}.

\begin{table}[h!]
\centering\small
\caption{The four accepted events in the Part slippage outline.}\label{tab:events}
\begin{tabularx}{\linewidth}{@{}r l l Y@{}}
\toprule
Event & LLM event name & Resource & Purpose \\
\midrule
1 & \code{stage_KET4_in_buffer} & \code{ur5e-3} & Stage the square peg in the buffer. \\
2 & \code{retract_ur5e4_home_clear_board} & \code{ur5e-4} & Retract to the home pose. \\
3 & \code{pick_dropped_gear_small_ur5e3} & \code{ur5e-3} & Pick the dropped gear. \\
4 & \code{place_gear_small_on_board_ur5e3} & \code{ur5e-3} & Place the gear in its assembly slot. \\
\bottomrule
\end{tabularx}
\end{table}

\Needspace{9\baselineskip}
\subsection{How feedback changed event 3}

On turn~4, the LLM proposed \code{pick_dropped_gear_small_ur5e3} as
\code{RECOVERY_SEQ3}. Its first candidate specified a robot pose in the
expected end state.

\excerptlabel{Reply}{Turn 4, first candidate}
\begin{lstlisting}
"pose": {"x": 0.25, "y": 0.45, "z": 1.243},
\end{lstlisting}

The RA refused both candidates on that turn. The next prompt carried the reason
for refusing the first candidate.

\excerptlabel{Prompt}{Turn 5, Validation Feedback}
\begin{lstlisting}
gear_small's recorded position gives ur5e-3's target, so this event takes no pose; it states {x: 0.250, y: 0.450, z: 1.243}
\end{lstlisting}

On turn~5, the LLM kept the event name and specified who would hold the
gear after the pick. It omitted the proposed pose.

\excerptlabel{Reply}{Turn 5, accepted candidate}
\begin{lstlisting}
"expected_end_state": {
  "resource_state": "picked",
  "held_part": "gear_small",
  "resource_location": null,
  "part_location": "ur5e-3"
},
\end{lstlisting}

The RA returned \code{FEASIBLE} and the CCA safety check returned \code{SAFE}.
After the fourth event placed the gear in the projected state, the LLM returned
\code{recovery_complete} on turn~7. The CCA check for continuation of normal
production returned \code{NONBLOCKING}.

\Needspace{14\baselineskip}
\subsection{Recorded outline prompt and reply format}\label{sec:outline-prompt}

These are the two messages from turn~1 of the event-level Framework run.
The user message gives the instructions and the observed state of the cell.
The system message specifies the JSON reply. Line wrapping, spacing and JSON
indentation have been adjusted for print. The text includes the complete
resource and part inventories from this turn.

\recordheading{User message \enspace{\normalfont\itshape Part slippage, outline turn 1}}
% begin recorded outline-user
\begin{lstlisting}[style=recorded]
Task and Role
You are the recovery planner. A failure interrupted normal production. Until normal
production takes over, only your recovery events change the planned state; a job
already running keeps running until an event pauses or stops it.
Plan how to get from the current state to the Recovery Goal. Propose 2(*@{\rmfamily\textendash}@*)5 different
ways to do the next step of that plan, most preferred first. The first one in your
order that passes every check becomes the next recovery event; the others are dropped.
Give your plan in `plan` (two short sentences) and keep each `rationale` to one.

Failure
What was detected and measured, stated the same way for every resource.
- ur5e-4 did not complete place_approach of gear_small (destination_location
assembly_board-v1)
- task_outcome (ur5e-4, task environment_9): goal_cancelled true, terminal_status 5,
progress 0.503
- device_fault_report (ur5e-4): resource_state failed, reported_by gripper object
detection
- state_discrepancy (ur5e-4, gear_small): expected_holder ur5e-4, measured_holder null
- gear_small: not held, not at any station, measured at (0.250, 0.450, 1.025)
- ur5e-4: measured pose (0.029, 0.145, 1.289), inside assembly_board-v1

Recovery Goal
- gear_small: part_location assembly_board-v1 (target Gear_Plate/Gear_Shaft_1)
(normally ur5e-4's place_insert; ur5e-4 failed)

Candidate Rules
- One event per candidate, by one resource: `outline_id`, `event_name` (what it
achieves; new names allowed), `resource`, `part_name` (if a part is involved),
`expected_end_state`, `neighbor_end_state` (if needed), `rationale` (the fact it
relies on).

- Use a resource only where it works (a reachable location, a position within its
workspace bounds, or a position it moves to) and for what it handles (parts it holds,
processes it runs). Working at one place says nothing about another.

- `expected_end_state`: the new value of each declared variable the event changes, for
its resource and its part only. Values may be new, but places, parts and positions
must be known. Give `pose` only for a position nothing else gives: an arm's pick or
place target comes from the part or the place when it declares compute_pick_targets or
compute_place_targets, a named pose gives its own position, and a mobile base's
reachable location gives its pose. An arm's `pose` is its controlled link's; its
fingertips lie the listed distance further along the tool. Write `pose` as {x, y, z},
a mobile base's as {x, y, yaw}.

- A value holds after an event only when it already held, a declared action of the
acting resource sets it, or the event's checks find it; state each value they find. A
pick empties what recorded the part where it was and leaves the part with the acting
resource; a place puts it where the event names, and a part put into or taken from a
place leaves that place's resource what its model declares; a move, navigation, pick
or place leaves the acting resource at that declared location, or at resource_location
null where no declared place holds it, as a stated `pose` does; a part placed where
its resource lists a completion gains that processCompleted entry; a stated completed
process entry ends that process (current_process null).

- `neighbor_end_state`: when moving the part changes another resource (it arrives at
or leaves it), that resource's own changed variables, not the part's: {"<resource
name>": {"<variable>": <value>}}. That resource does not act; variables you leave out
stay as they are.

- Use variable names exactly as declared. Look up any fact or variable you rely on
that is not shown.

- When you judge that normal production can carry on its remaining work from the state
your events leave, return only `plan` and `decision: "recovery_complete"`; normal
production then takes over.

Knowledge Base
Request facts from the knowledge base by returning only `plan`, `decision:
"need_context"` and 1(*@{\rmfamily\textendash}@*)8 `context_requests` from the list below; they are added to the
next prompt before you propose events. When proposing events, omit `decision` and
`context_requests`.

- /resources/<name>/snapshot: every recorded value of a resource; once your events
change it, its expected values after them

- /resources/<name>/state_variables: its condition variables, their allowed values,
and current values

- /resources/<name>/pending_tasks: its remaining planned tasks

- /resources/<name>/static_capabilities: named poses, staging areas and process
options

- /resources/<name>/pose_relations: which areas contain its measured position (one
point, not its full extent), which declared places contain it, and how far it lies
from each resource's workspace bounds

- /parts/<name>/snapshot: every recorded value of a part

- /parts/<name>/state_variables: its condition variables and current values

- /parts/<name>/target: where the part belongs in the finished product (its assembly
location and slot pose)

- /parts/<name>/pose_relations: which areas contain its measured position (one point,
not its full extent), which declared places contain it, and how far it lies from each
resource's workspace bounds

- /regions: area bounds in world coordinates

- /geometry: station poses and sizes, robot and mobile bases, floor extent and floor
obstacles (a mobile base's footprint keeps a 0.03 m margin from them)

Resources
Recorded at the failure. Every resource is listed the same way: its state, then the
capabilities it declares (reachable locations, where it places parts, workspace
bounds, positions it moves to, processes it runs, job control, its actions with what
each needs and sets). A resource's condition is the variables its model declares;
request `/resources/<name>/state_variables` for them and their allowed values. A value
your events changed without stating it reads "not stated".

3D Printing Station (resource) completed; output.<part name>: gear_medium, gear_large;
a part taken out of 3D Printing Station: from completed with another output left it
ends completed; from completed with no other output left it ends idle

Buffer For Machined parts (resource) idle; zone_1_part null; zone_2_part null;
zone_3_part null; zone_4_part null; capabilities: moves parts in one direction:
zone_1_part (*@$\rightarrow$@*) zone_2_part (*@$\rightarrow$@*) zone_3_part (*@$\rightarrow$@*) zone_4_part; actions: detect_parts;
compute_transport_targets; move_relative

Conveyor (resource) idle; part_location.<part name>: none (values loading_position_1,
after loading_position_1, loading_position_2, after loading_position_2, output_nest);
part_order.<part name>: none; capabilities: moves parts in one direction:
loading_position_1 (*@$\rightarrow$@*) after loading_position_1 (*@$\rightarrow$@*) loading_position_2 (*@$\rightarrow$@*) after
loading_position_2 (*@$\rightarrow$@*) output_nest; part_order numbers the parts on it from 0, the one
furthest along first; a part put on it is numbered after the parts ahead of it, and
the parts behind it move up by one; actions: detect_parts; compute_transport_targets;
move_relative

Exit (resource) empty; part_name null; a part put into Exit: from empty it ends
occupied

KMR (kmr) idle; held_part null; resource_location Storage; capabilities: reachable
locations Storage, M1, M2; moves to any position (x, y, yaw); actions:
compute_pick_targets; compute_place_targets; detect_parts; grasp_part (needs held_part
null) sets resource_state carrying, held_part <part_name>; move_base sets pose (x, y,
yaw); move_cartesian; move_joints; move_relative; move_to_named_pose (needs held_part
null) sets resource_state idle, or (needs held_part not null) sets resource_state
carrying; release_part (needs held_part not null) sets resource_state placed,
held_part null; rotate_joint

M1 (resource) idle; part_name null; staging_part null (M1 staging tray); capabilities:
trim (circle, square); job control: a paused job continues from where it was paused,
also after its part is unloaded and loaded back in; a stopped job cannot be restarted;
actions: select_program (needs resource_state not running) sets selected_program
<program>; start_job (needs resource_state loaded, part_name <part_name>,
selected_program circle, the part's current_process not matching {machine: M1}) sets
resource_state running, the part's current_process {process: <process>, result:
<result>, status: running, machine: M1}, or (needs resource_state loaded, part_name
<part_name>, selected_program square, the part's current_process not matching
{machine: M1}) sets resource_state running, the part's current_process {process:
<process>, result: <result>, status: running, machine: M1}; pause_job (needs
resource_state running, part_name <part_name>, the part's current_process matching
{machine: M1}) sets resource_state paused, the part's current_process status paused;
resume_job (needs resource_state paused, part_name <part_name>, selected_program
circle, the part's current_process matching {process: trim, result: circle, machine:
M1}) sets resource_state running, the part's current_process status running, or (needs
resource_state paused, part_name <part_name>, selected_program square, the part's
current_process matching {process: trim, result: square, machine: M1}) sets
resource_state running, the part's current_process status running, or (needs
resource_state loaded, part_name <part_name>, selected_program circle, the part's
current_process matching {status: paused, process: trim, result: circle, machine: M1})
sets resource_state running, the part's current_process status running, or (needs
resource_state loaded, part_name <part_name>, selected_program square, the part's
current_process matching {status: paused, process: trim, result: square, machine: M1})
sets resource_state running, the part's current_process status running; stop_job
(needs resource_state running or paused, part_name <part_name>, the part's
current_process matching {machine: M1}) sets resource_state stopped, the part's
current_process status stopped; complete_job uncontrollable (needs resource_state
running, part_name <part_name>, the part's current_process matching {machine: M1})
sets resource_state completed, the part's current_process null, the part's
processCompleted gains {process: trim, result: <result>}; a part put into M1: from
idle it ends loaded; a part taken out of M1: from completed or paused or stopped or
loaded it ends idle; from interrupted it keeps its state

M2 (resource) idle; part_name null; staging_part null (M2 staging tray); capabilities:
trim (circle, square); job control: a paused job continues from where it was paused,
also after its part is unloaded and loaded back in; a stopped job cannot be restarted;
actions: select_program (needs resource_state not running) sets selected_program
<program>; start_job (needs resource_state loaded, part_name <part_name>,
selected_program circle, the part's current_process not matching {machine: M2}) sets
resource_state running, the part's current_process {process: <process>, result:
<result>, status: running, machine: M2}, or (needs resource_state loaded, part_name
<part_name>, selected_program square, the part's current_process not matching
{machine: M2}) sets resource_state running, the part's current_process {process:
<process>, result: <result>, status: running, machine: M2}; pause_job (needs
resource_state running, part_name <part_name>, the part's current_process matching
{machine: M2}) sets resource_state paused, the part's current_process status paused;
resume_job (needs resource_state paused, part_name <part_name>, selected_program
circle, the part's current_process matching {process: trim, result: circle, machine:
M2}) sets resource_state running, the part's current_process status running, or (needs
resource_state paused, part_name <part_name>, selected_program square, the part's
current_process matching {process: trim, result: square, machine: M2}) sets
resource_state running, the part's current_process status running, or (needs
resource_state loaded, part_name <part_name>, selected_program circle, the part's
current_process matching {status: paused, process: trim, result: circle, machine: M2})
sets resource_state running, the part's current_process status running, or (needs
resource_state loaded, part_name <part_name>, selected_program square, the part's
current_process matching {status: paused, process: trim, result: square, machine: M2})
sets resource_state running, the part's current_process status running; stop_job
(needs resource_state running or paused, part_name <part_name>, the part's
current_process matching {machine: M2}) sets resource_state stopped, the part's
current_process status stopped; complete_job uncontrollable (needs resource_state
running, part_name <part_name>, the part's current_process matching {machine: M2})
sets resource_state completed, the part's current_process null, the part's
processCompleted gains {process: trim, result: <result>}; a part put into M2: from
idle it ends loaded; a part taken out of M2: from completed or paused or stopped or
loaded it ends idle; from interrupted it keeps its state

Storage (resource) idle; inventory.<part name>: KET8_Square_8mm, KET12_Square_12mm,
KET16_Square_16mm, RGOCG4-50_Round_4mm, RGOCG8-50_8mm, RGOCG12-50_12mm,
RGOCG16-50_16mm

ur5e-1 (robot) idle; held_part null; resource_location home; pose (-4.348, 0.699,
1.280); capabilities: reachable locations M1, M1 staging tray, Conveyor; places parts
on Conveyor at loading_position_1; workspace bounds; actions: compute_pick_targets;
compute_place_targets; detect_parts; grasp_part (needs held_part null) sets
resource_state picked, held_part <part_name>; move_cartesian sets pose (x, y, z);
move_joints; move_relative (needs pose recorded) sets pose moved by (dx, dy, dz);
move_to_named_pose (needs held_part null) sets resource_state idle, resource_location
<pose_name>; release_part (needs held_part not null) sets resource_state placed,
held_part null; rotate_joint; fingertips 0.218 m beyond its pose along the tool

ur5e-2 (robot) idle; held_part null; resource_location home; pose (-2.348, 0.699,
1.280); capabilities: reachable locations M2, M2 staging tray, Conveyor; places parts
on Conveyor at loading_position_2; workspace bounds; actions: compute_pick_targets;
compute_place_targets; detect_parts; grasp_part (needs held_part null) sets
resource_state picked, held_part <part_name>; move_cartesian sets pose (x, y, z);
move_joints; move_relative (needs pose recorded) sets pose moved by (dx, dy, dz);
move_to_named_pose (needs held_part null) sets resource_state idle, resource_location
<pose_name>; release_part (needs held_part not null) sets resource_state placed,
held_part null; rotate_joint; fingertips 0.218 m beyond its pose along the tool

ur5e-3 (robot) picked; held_part KET4_Square_4mm; resource_location Buffer For
Machined parts; pose (-0.320, 0.500, 1.501); capabilities: reachable locations Buffer
For Machined parts, assembly_board-v1, Exit; workspace bounds; actions:
compute_pick_targets; compute_place_targets; detect_parts; grasp_part (needs held_part
null) sets resource_state picked, held_part <part_name>; move_cartesian sets pose (x,
y, z); move_joints; move_relative (needs pose recorded) sets pose moved by (dx, dy,
dz); move_to_named_pose (needs held_part null) sets resource_state idle,
resource_location <pose_name>; release_part (needs held_part not null) sets
resource_state placed, held_part null; rotate_joint; a part it places at
assembly_board-v1 gains processCompleted {process: assembly, target: <its assembly
target>}; fingertips 0.218 m beyond its pose along the tool

ur5e-4 (robot) failed; held_part null; resource_location null; pose (0.029, 0.145,
1.289), inside assembly_board-v1; capabilities: reachable locations 3D Printing
Station, assembly_board-v1; workspace bounds; actions: compute_pick_targets;
compute_place_targets; detect_parts; grasp_part (needs held_part null) sets
resource_state picked, held_part <part_name>; move_cartesian sets pose (x, y, z);
move_joints; move_relative (needs pose recorded) sets pose moved by (dx, dy, dz);
move_to_named_pose (needs held_part null) sets resource_state idle, resource_location
<pose_name>; release_part (needs held_part not null) sets resource_state placed,
held_part null; rotate_joint; a part it places at assembly_board-v1 gains
processCompleted {process: assembly, target: <its assembly target>}; fingertips 0.218
m beyond its pose along the tool

Parts
A part's condition is its variables: where it is, who holds it, and its process record
(its unfinished process and the processes it completed); request
`/parts/<name>/state_variables` for them.

KET12_Square_12mm at Storage, assembly target
GMC_Laser_Plate_Virtual/KET12_Square_12mm

KET16_Square_16mm at Storage, assembly target
GMC_Laser_Plate_Virtual/KET16_Square_16mm

KET4_Square_4mm held by ur5e-3, assembly target
GMC_Laser_Plate_Virtual/KET4_Square_4mm; processes: trim (*@$\rightarrow$@*) square completed

KET8_Square_8mm  at Storage, assembly target GMC_Laser_Plate_Virtual/KET8_Square_8mm

RGOCG12-50_12mm  at Storage, assembly target GMC_Laser_Plate_Virtual/RGOCG12-50_12mm

RGOCG16-50_16mm  at Storage, assembly target GMC_Laser_Plate_Virtual/RGOCG16-50_16mm

RGOCG4-50_Round_4mm at Storage, assembly target
GMC_Laser_Plate_Virtual/RGOCG4-50_Round_4mm

RGOCG8-50_8mm  at Storage, assembly target GMC_Laser_Plate_Virtual/RGOCG8-50_8mm

assembly_board-v1  at assembly_board-v1

gear_large at 3D Printing Station, assembly target Gear_Plate/Gear_Shaft_3; processes:
print_part completed

gear_medium at 3D Printing Station, assembly target Gear_Plate/Gear_Shaft_2;
processes: print_part completed

gear_small not held, not at any station, measured at (0.250, 0.450, 1.025), assembly
target Gear_Plate/Gear_Shaft_1; processes: print_part completed
\end{lstlisting}
% end recorded outline-user

\recordheading{System message \enspace{\normalfont\itshape Outline reply format}}
% begin recorded outline-system
\begin{lstlisting}[style=recorded]
Reply with only one JSON object that follows this JSON schema:
{
  "type": "object",
  "additionalProperties": false,
  "$defs": {
    "outline_event": {
      "type": "object",
      "additionalProperties": false,
      "properties": {
        "outline_id": {"type": "string"},
        "event_name": {"type": "string", "minLength": 1},
        "resource": {"type": "string"},
        "part_name": {"type": ["string", "null"], "minLength": 1},
        "expected_end_state": {"type": "object", "minProperties": 1},
        "neighbor_end_state": {
          "type": "object",
          "additionalProperties": {"type": "object"},
          "description": "Changes to other resources' variables caused by moving the
          event's part, by resource name: {\"<resource name>\": {\"<variable>\":
          <value>}}. Those resources only record the change; the event's resource is
          the one that acts."
        },
        "rationale": {"type": "string"}
      },
      "required": ["outline_id", "event_name", "resource", "expected_end_state",
      "rationale"]
    }
  },
  "properties": {
    "plan": {"type": "string"},
    "candidate_events": {
      "type": "array",
      "minItems": 2,
      "maxItems": 5,
      "items": {"$ref": "#/$defs/outline_event"}
    },
    "decision": {
      "type": ["string", "null"],
      "enum": [null, "need_context", "recovery_complete"],
      "description": "Use null for proposed events, need_context to request facts, or
      recovery_complete when you judge the failure overcome."
    },
    "context_requests": {
      "type": "array",
      "minItems": 0,
      "maxItems": 8,
      "description": "Use [] for proposed events; need_context requires 1 to 8
      references and no candidate fields.",
      "items": {"type": "string", "pattern": "^/"}
    }
  },
  "required": ["plan"]
}
\end{lstlisting}
% end recorded outline-system

\subsection{Prompt differences between settings}\label{sec:differences}

At the event level, the framework considers the candidates in the LLM's order
until one passes its checks. Baseline and Informed baseline take the first
candidate. The event-level baseline prompts replace the framework's selection
instruction with the following sentence.

\excerptlabel{Prompt}{Baseline and Informed baseline, event-level}
\begin{lstlisting}
The first one becomes the next recovery event; the others are dropped.
\end{lstlisting}

Informed baseline also inserts the following section after the recovery goal.

% begin recorded informed-safety
\begin{lstlisting}[style=recorded]
Safety Rules
Your events, and normal production after them, must keep these rules.
- SAFE_shared_area_mutex: At most one resource may occupy the assembly_board-v1
destination area at any time.
- SAFE_M1_mutex: At most one resource accessing M1 may occupy its protected area at
any time.
- SAFE_M2_mutex: At most one resource accessing M2 may occupy its protected area at
any time.
- SAFE_trim_before_assembly_board-v1_entry: A part with a trim step may enter
assembly_board-v1 only after its trim with the required result is completed.
- SAFE_gear_small_before_KET4_Square_4mm: KET4_Square_4mm may complete assembly only
after assembly of gear_small is completed.
- SAFE_KET4_Square_4mm_before_RGOCG4-50_Round_4mm: RGOCG4-50_Round_4mm may complete
assembly only after assembly of KET4_Square_4mm is completed.
- SAFE_KET4_Square_4mm_before_RGOCG4-50_Round_4mm_processing: RGOCG4-50_Round_4mm must
not complete trim with result circle until KET4_Square_4mm has completed trim with
result square.
\end{lstlisting}
% end recorded informed-safety

At the outline level, all three methods request the complete sequence in
\code{events}. The final event rule defines where that sequence should end.

\excerptlabel{Prompt}{Outline-level, final event rule}
\begin{lstlisting}
- End `events` where you judge that normal production can carry on its remaining work from the state your events leave; normal production then takes over.
\end{lstlisting}

The same event definition is used at both levels. The outline-level reply
format replaces \code{candidate_events} with \code{events} and allows
\code{decision} to be \code{null} or \code{need_context}. Submitting the
ordered events supplies the outline for checking.


\recordheading{System message excerpt \enspace{\normalfont\itshape Outline-level reply properties}}
% begin recorded outline-level-properties
\begin{lstlisting}[style=recorded]
"properties": {
  "plan": {"type": "string"},
  "events": {"type": "array", "minItems": 1, "items": {"$ref":
    "#/$defs/outline_event"}, "description": "The whole recovery outline: every
    recovery event, in the order they happen."},
  "decision": {"type": ["string", "null"], "enum": [null, "need_context"],
    "description": "Use null for proposed events; use need_context only to request
    facts without proposing events."},
  "context_requests": {"type": "array", "minItems": 0, "maxItems": 8, "description":
    "Use [] for proposed events; need_context requires 1 to 8 references and no
    events.", "items": {"type": "string", "pattern": "^/"}}
}
\end{lstlisting}
% end recorded outline-level-properties

\clearpage
\section{Primitive Composition}
\setcounter{table}{0}
\renewcommand*{\theHtable}{B.\arabic{table}}
\subsection{The accepted event becomes the next prompt}\label{sec:composition-case}

Event~3, \code{pick_dropped_gear_small_ur5e3} (\code{RECOVERY_SEQ3}), uses
\code{ur5e-3}'s primitives. The end state is \code{picked}, with
\code{held_part} set to \code{gear_small}, \code{resource_location} set to
\code{null}, and \code{part_location} set to the robot's identifier,
\code{recovery-resource-3@localhost}.

\excerptlabel{Prompt}{Instruction for the active event}
\begin{lstlisting}
Author primitive_steps for ONLY the active DES transition, using only the visible primitive catalog for the active resource.
\end{lstlisting}

Composition calls~5--7 all concern this pick. On call~5 the LLM asked for the
primitive contracts, the robot's nominal pick decompositions, its projected
state, and the gear observation. The request included these references.

\excerptlabel{Reply}{Call 5, first three primitive contracts requested}
\begin{lstlisting}
"/primitive_contracts/compute_pick_targets",
"/primitive_contracts/move_cartesian",
"/primitive_contracts/grasp_part",
\end{lstlisting}

The served contract describes what the target calculation does. The moves and
the grasp still have to be specified as separate primitive steps.

\excerptlabel{Prompt}{Call 6, served contract for \texttt{compute\_pick\_targets}}
\begin{lstlisting}
Compute vertical pick target positions from world object geometry and configured end-effector/TCP feedback without moving. The frames must agree.
\end{lstlisting}

\subsection{The revised and accepted program}

On call~6, the RA flagged two inputs to \code{compute_pick_targets}:
\code{target_pose} and \code{target_pose_source}. The LLM had bound the coordinates to the gear's observed pose and explicitly
named that source. The RA checked those parameters against the inputs used by
the calculation. Its finding reached the LLM on call~7.

\excerptlabel{Reply}{Call 6, bindings supplied to the target calculation}
\begin{lstlisting}
"target_pose": {
  "x": {"context_ref": "/parts/gear_small/observed_pose/x"},
  "y": {"context_ref": "/parts/gear_small/observed_pose/y"},
  "z": {"context_ref": "/parts/gear_small/observed_pose/z"},
  "model_name": {"context_ref": "/parts/gear_small/target/model_name"}
},
"target_pose_source": "observed_pose"
\end{lstlisting}

\excerptlabel{Prompt}{Call 7, feedback on call 6}
\begin{lstlisting}
compute_pick_targets: target_pose_source 'observed_pose' differs from '', the input the composition's calculation reads
\end{lstlisting}

\Needspace{5\baselineskip}
With only the part name supplied, the calculation reads the recorded gear pose.

\excerptlabel{Reply}{Call 7, first step}
\begin{lstlisting}
{"primitive": "compute_pick_targets", "params": {"part_name": "gear_small"}},
\end{lstlisting}

\begin{table}[h!]
\centering\small
\caption{Accepted program for event 3, \code{pick_dropped_gear_small_ur5e3}.}\label{tab:pick}
\begin{tabularx}{\linewidth}{@{}r l Y@{}}
\toprule
Step & Primitive & What the robot does \\
\midrule
1 & \code{compute_pick_targets} & Calculate approach, target and retreat poses for \code{gear_small}. \\
2 & \code{move_cartesian} & Move above the gear to \code{approach_pose}. \\
3 & \code{move_cartesian} & Descend to \code{target_pose}. \\
4 & \code{grasp_part} & Grasp \code{gear_small}. \\
5 & \code{move_cartesian} & Withdraw with the gear to \code{retreat_pose}. \\
\bottomrule
\end{tabularx}
\end{table}

\excerptlabel{Reply}{Call 7, step 2 binds \texttt{x} to the approach pose calculated in step 1}
\begin{lstlisting}
"x": {"context_ref": "event_facts.pick_targets.gear_small.approach_pose.x"},
\end{lstlisting}

\Needspace{13\baselineskip}
\subsection{Recorded composition prompt and reply format}\label{sec:composition-prompt}

Call~5 is the first composition call for event~3,
\code{pick_dropped_gear_small_ur5e3}. The user message first lists the
primitives and references available to \code{ur5e-3}, then gives the transition
and its start and end states. Its output contract specifies how to return a
program or request more context. Both messages are reproduced in full with
the same print formatting as the outline prompt.

\recordheading{User message \enspace{\normalfont\itshape Part slippage, composition call 5}}
% begin recorded composition-user
\begin{lstlisting}[style=recorded]
Task and Role
You are the active replanner authoring primitive_steps for a DES fallback recovery
session.
Current phase: Event-to-Primitive Authoring (event-local).
Author primitive_steps for ONLY the active DES transition, using only the visible
primitive catalog for the active resource. There is no mechanical composer behind you:
if you emit no primitive_steps, no trace is produced for this event.
RA primitive_support records independent capability witnesses. Primitive Composition
must still establish ordering, compatible parameter bindings, and achievement of the
complete expected_end_state. An accepted outline is not an executable program.
The active event is provided inline below. For other facts, emit decision=need_context
with refs in context_requests; they are resolved in the next turn under 'Served
Context'. Do not guess values. Capability decompositions show how this resource's
nominal functions compose primitives. The <function_name> token must be one of the
names listed under 'Available Capability Decompositions'; primitive names such as
grasp_part or release_part belong under /primitive_contracts/<name>.

Visible Primitive Catalog (names only; full contract via /primitive_contracts/<name>)
[
  {"name": "compute_pick_targets"},
  {"name": "compute_place_targets"},
  {"name": "detect_parts"},
  {"name": "grasp_part"},
  {"name": "move_cartesian"},
  {"name": "move_relative"},
  {"name": "move_to_named_pose"},
  {"name": "release_part"}
]

Published Ref Schema (for context_requests)
[
  {"ref": "/outline_event", "description": "the active outline event (full)"},
  {"ref": "/resources/<jid>/snapshot", "description": "full resource snapshot"},
  {
    "ref": "/resources/<jid>/current_pose",
    "description": "an arm's current tool pose; a KMR's poses are in its snapshot"
  },
  {"ref": "/resources/<jid>/held_part", "description": "currently held part, if any"},
  {
    "ref": "/resources/<jid>/static_capabilities",
    "description": "as the capture recorded: the places an arm reaches, a machine's
    process capabilities and configuration; empty for other resources; no poses or
    coordinates"
  },
  {
    "ref": "/resources/<jid>/resource_evidence/controller_config/cartesian_motion",
    "description": "an arm's declared route around its base, as x, y, z points:
    pick_transit_waypoints and transit_waypoints, and rotation_waypoint, where it
    turns"
  },
  {"ref": "/parts/<name>/snapshot", "description": "full part state"},
  {
    "ref": "/parts/<name>/observed_pose",
    "description": "last observed pose for the part, if any"
  },
  {"ref": "/parts/<name>/target", "description": "destination geometry if grounded"},
  {
    "ref": "/primitive_catalog/<jid>",
    "description": "visible primitive catalog for the resource"
  },
  {
    "ref": "/primitive_contracts/<name>",
    "description": "contract card for one visible primitive"
  },
  {
    "ref": "/projected_snapshot/<jid>",
    "description": "this resource at the start of the active event"
  },
  {
    "ref": "/capability_decompositions/<function_name>",
    "description": "how a nominal function of the active resource composes primitives"
  },
  {
    "ref": "/regions/<name>",
    "description": "region bounds in the world frame (or /regions for all)"
  },
  {
    "ref": "/geometry/<key>",
    "description": "station poses and sizes, robot and mobile bases, floor extent and
    obstacles (or /geometry for all)"
  }
]

Available Capability Decompositions (names only)
["move_home", "pick_approach", "pick_grasp", "place_approach", "place_insert"]

Active DES Transition (token)
{
  "outline_id": "RECOVERY_SEQ3",
  "resource_jid": "recovery-resource-3@localhost",
  "action": "pick_dropped_gear_small_ur5e3",
  "part_name": "gear_small",
  "start_state_id": null,
  "end_state_id": null
}

Active DES Transition (full event)
{
  "outline_id": "RECOVERY_SEQ3",
  "resource_jid": "recovery-resource-3@localhost",
  "event_name": "pick_dropped_gear_small_ur5e3",
  "description": "gear_small's measured position is 0.255 m from ur5e-3's base in xy,
  at z 1.025, which is inside its envelope (radius 0.886 m, z 0.971(*@{\rmfamily\textendash}@*)1.121); the pick
  target comes from the part's recorded position, so no pose is given.",
  "part_name": "gear_small",
  "target_ref": "recovery-resource-3@localhost",
  "expected_start_state": {
    "held_part": null,
    "part_location": null,
    "resource_location": "Buffer For Machined parts",
    "resource_state": "placed"
  },
  "expected_end_state": {
    "resource_state": "picked",
    "held_part": "gear_small",
    "resource_location": null,
    "part_location": "recovery-resource-3@localhost"
  }
}

Transition Authoring Cursor
{"active_index": 0, "accepted_outline_count": 1, "remaining_events_count": 1}

Accepted Outline Summary (context for active event)
- RECOVERY_SEQ3 / recovery-resource-3@localhost / gear_small ->
pick_dropped_gear_small_ur5e3 to recovery-resource-3@localhost
  (gear_small's measured position is 0.255 m from ur5e-3's base in xy, at z 1.025,
  which is inside its envelope (radius 0.886 m, z 0.971(*@{\rmfamily\textendash}@*)1.121); the pick target comes
  from the part's recorded position, so no pose is given.)

Named Pose Rules
[
  "home"
]
- move_to_named_pose may only use named poses advertised for the active resource.

Output Contract
- Keep `plan` to at most two short sentences and `rationale` to one short sentence.
- Author primitive_steps for ONLY the active DES transition shown above.
- outline_id and resource_jid MUST exactly match the active transition.
- primitive_steps[*].primitive MUST be a name in Visible Primitive Catalog.
- Each step MUST include params. Every number in params is a {"context_ref": ...};
strings and booleans may be literals from served context.
- compute_pick_targets publishes event_facts.pick_targets.<part>;
compute_place_targets publishes event_facts.place_targets.<part>.
- need_context: refs in context_requests, primitive_steps empty.
- primitive_steps_ready: non-empty primitive_steps. need_primitive_revision: say why
in rationale.

RA composition validation: author an ordered program that establishes every intended
resource and part postcondition from declared effects or what the checks find. A
robot's resource_location is the last place it reached: a transfer's declared effect,
a named pose's declared place, or the place whose recorded area holds the gripper
where a motion ends; a motion that ends in no recorded area leaves it as it was. A
KMR's is the place whose declared dock its base ends on, or none. Primitive support is
only a union of possible contributors, not a sequence or a requirement to use every
primitive. Use exact primitive identities and explicit {"context_ref": "/path"}
bindings. Every number in params is such a reference, and a number written into a
program is refused; strings and booleans may be literals. The numbers of a resource's
nominal steps are in its model, under
/resources/<jid>/recovery_des_model/events/<index>/composition_contract/steps/<index>/params.
A contract with controllable false is the resource's own outcome, after after_sec: a
step that names it sends no command and waits for it. The recovery_des_model is
retrievable under /resources/recovery-resource-3@localhost/recovery_des_model.
Recorded owner inputs are under
/resources/recovery-resource-3@localhost/resource_evidence/<field>: request only the
fields a parameter binds to; the whole record is too large and is listed by field.
Missing facts require need_context; they are different from explicit null. Calculated
event_facts are available only after their producing step. The owner's target
calculations read recorded geometry and observed poses themselves. An accepted
projected placement can support compute_pick_targets in this model-based composition.
The robot target outputs include approach_pose, target_pose, and retreat_pose. Only a
place into the part's assembly_slot (the location in /parts/<name>/target) has slot
fields in the calculated place targets. Validation findings identify the step or
missing postcondition to revise. This is composition for review; no execution is
authorized.
\end{lstlisting}
% end recorded composition-user

\recordheading{System message \enspace{\normalfont\itshape Composition reply format}}
% begin recorded composition-system
\begin{lstlisting}[style=recorded]
Reply with only one JSON object that follows this JSON schema:
{
  "type": "object",
  "properties": {
    "plan": {"type": "string"},
    "decision": {
      "type": "string",
      "enum": ["primitive_steps_ready", "need_context", "need_primitive_revision"]
    },
    "outline_id": {"type": "string"},
    "resource_jid": {"type": "string"},
    "context_requests": {"type": "array", "items": {"type": "string"}},
    "rationale": {"type": "string"},
    "notes": {"type": "array", "items": {"type": "string"}},
    "primitive_steps": {
      "type": "array",
      "items": {
        "type": "object",
        "properties": {"primitive": {"type": "string"}, "params": {"type": "object"}},
        "required": ["primitive", "params"]
      }
    }
  },
  "required": ["plan", "decision", "outline_id", "resource_jid"]
}
\end{lstlisting}
% end recorded composition-system

\Needspace{24\baselineskip}
\section{Recovery Runs Executed in Gazebo}
\setcounter{table}{0}
\renewcommand*{\theHtable}{C.\arabic{table}}

For each scenario, we compared the passing event-level Framework outlines from
all three models and selected one with the fewest events.
The table summarizes the synthesis and composition of those outlines and the
number of events completed during replay in Gazebo. The event names below
are those proposed by the model.

\begin{center}
\small
\begin{tabularx}{\linewidth}{@{}Y l c c c@{}}
\toprule
Scenario & Model & \shortstack{Outline\\turns} & \shortstack{Composition\\calls} & \shortstack{Events\\completed} \\
\midrule
Part slippage & Claude Opus 5.5 & 7 & 9 & 4 of 4 \\
Conveyor breakdown & Claude Opus 5.5 & 12 & 15 & 7 of 7 \\
ur5e-1 breakdown & GPT-6.1 Sol & 11 & 19 & 4 of 4 \\
Machining breakdown during part processing & GPT-6.1 Sol & 20 & 49 & 14 of 14 \\
\bottomrule
\end{tabularx}
\end{center}

\subsection{Part slippage}

\code{ur5e-4} dropped \code{gear_small} on its approach to the assembly board.
To recover it, \code{ur5e-3} first put down the square peg it was holding.
The live placement calculation could not resolve the buffer geometry, and the
staging motion used the saved parameters. Table~\ref{tab:pick} gives the
five primitive steps used for event~3.

\smallskip
{\small
\renewcommand{\arraystretch}{1.12}
\begin{tabularx}{\linewidth}{@{}r Y l Y@{}}
\toprule
Event & Accepted event name & Resource & What happened in Gazebo \\
\midrule
1 & \code{stage_KET4_in_buffer} & \code{ur5e-3} & Released \code{KET4_Square_4mm} in the buffer, freeing its gripper. \\
2 & \code{retract_ur5e4_home_clear_board} & \code{ur5e-4} & Moved to its \code{home} pose, clearing the assembly board. \\
3 & \code{pick_dropped_gear_small_ur5e3} & \code{ur5e-3} & Grasped \code{gear_small} at its measured position and withdrew with it attached. \\
4 & \code{place_gear_small_on_board_ur5e3} & \code{ur5e-3} & Released the gear at \code{Gear_Plate/Gear_Shaft_1} and retreated. \\
\bottomrule
\end{tabularx}\par}

\clearpage
\subsection{Conveyor breakdown}

The Conveyor became unavailable while \code{ur5e-1} held the machined
square peg, \code{KET4_Square_4mm}. The M1 staging tray provided a handoff
point. KMR collected the peg there and carried it directly to
\texttt{Buffer For Machined parts}.

\smallskip
{\small
\renewcommand{\arraystretch}{1.12}
\begin{tabularx}{\linewidth}{@{}r Y l Y@{}}
\toprule
Event & Accepted event name & Resource & What happened in Gazebo \\
\midrule
1 & \code{place_on_M1_staging_tray} & \code{ur5e-1} & Released the square peg on the M1 staging tray. \\
2 & \code{navigate_KMR_to_M1} & \code{KMR} & Drove to M1 to collect the staged peg. \\
3 & \code{ur5e1_retract_home} & \code{ur5e-1} & Retracted to \code{home}, clearing access to the tray. \\
4 & \code{kmr_pick_from_M1_staging_tray} & \code{KMR} & Grasped the peg on the tray and lifted it clear. \\
5 & \code{kmr_move_base_to_buffer} & \code{KMR} & Carried the peg from M1 to the buffer. \\
6 & \code{kmr_place_on_buffer_zone_1} & \code{KMR} & Released the peg in \code{zone_1} of the buffer. \\
7 & \code{kmr_stow_arm_and_return_to_storage} & \code{KMR} & Stowed its arm in the \code{transport} pose and drove back to Storage. \\
\bottomrule
\end{tabularx}\par}

\subsection{ur5e-1 breakdown}

The gripper on \code{ur5e-1} failed before it could transfer the completed
square peg from M1 to the Conveyor. The empty arm could still retract.
KMR took over the transfer and delivered \code{KET4_Square_4mm} to the
original loading position.

\smallskip
{\small
\renewcommand{\arraystretch}{1.12}
\begin{tabularx}{\linewidth}{@{}r Y l Y@{}}
\toprule
Event & Accepted event name & Resource & What happened in Gazebo \\
\midrule
1 & Retract the empty arm to its named home pose & \code{ur5e-1} & Moved the empty arm to \code{home}, clearing access to M1. \\
2 & Retrieve the completed part from M1 with KMR & \code{KMR} & Drove to M1, grasped the square peg and stowed its arm for transport. \\
3 & Deliver to conveyor loading position 1 using a south-facing base heading & \code{KMR} & Approached the Conveyor facing south and released the peg at loading position~1. \\
4 & Return KMR to Storage and verify idle readiness & \code{KMR} & Stowed its arm and drove back to Storage. \\
\bottomrule
\end{tabularx}\par}

\clearpage
\subsection{Machining breakdown during part processing}

M1 failed during the square trim of \code{KET4_Square_4mm}. M2 could run
the square program, but it was trimming \code{RGOCG4-50_Round_4mm}.
The recovery temporarily unloaded the round peg, finished the square peg on M2,
then reloaded the round peg and resumed its interrupted job.

\smallskip
{\small
\renewcommand{\arraystretch}{1.12}
\begin{tabularx}{\linewidth}{@{}r Y l Y@{}}
\toprule
Event & Accepted event name & Resource & What happened in Gazebo \\
\midrule
1 & Pause M2's round-part job for temporary unloading & \code{M2} & Paused the running trim job for the round peg. \\
2 & Unload the paused round part into M2's staging tray & \code{ur5e-2} & Lifted the round peg out of M2 and placed it on M2's staging tray. \\
3 & Navigate to M1 and retrieve the interrupted square part & \code{KMR} & Drove to M1 and grasped the partially processed square peg. \\
4 & Select M2's square program before loading & \code{M2} & Switched the selected program from \code{circle} to \code{square}. \\
5 & Return ur5e-2 home to clear M2 access & \code{ur5e-2} & Moved to \code{home}, clearing the approach to M2. \\
6 & Navigate to M2 and load the interrupted square part & \code{KMR} & Carried the square peg to M2 and released it into the machine. \\
7 & Navigate KMR to Storage to clear M2 & \code{KMR} & Drove back to Storage, clearing the machine access area. \\
8 & Start the square-part job on M2 & \code{M2} & Started the \code{trim} job with the required \code{square} result. \\
9 & Await and verify completion of square trimming on M2 & \code{M2} & Recorded completion of the square trim on M2. \\
10 & Transfer the completed square part from M2 to Conveyor & \code{ur5e-2} & Unloaded the finished square peg and placed it on the Conveyor. \\
11 & Reload the paused round part from M2's staging tray & \code{ur5e-2} & Picked up the round peg from the staging tray and loaded it back into M2. \\
12 & Return ur5e-2 home to clear M2 for resumed processing & \code{ur5e-2} & Withdrew to \code{home}, clearing the machine for processing. \\
13 & Restore M2's circle program before a separate resume event & \code{M2} & Selected \code{circle} for the reloaded round peg. \\
14 & Resume the reloaded paused round-part job & \code{M2} & Resumed the same round-part job that was paused in event~1. \\
\bottomrule
\end{tabularx}\par}
\end{document}
